\subsubsection*{Численный пример}

\begin{idea}[title={Откуда этот материал}]
Прогон метода условного градиента на числах; на лекции примера не было.
Допустимое множество взято многогранным, чтобы вспомогательная задача была
обычной задачей ЛП.
\end{idea}

\begin{example}\label{ex:frank-wolfe-run}
\[
  \varphi(x) = (x_1 - 0{,}3)^2 + (x_2 - 0{,}7)^2 \ \longrightarrow\ \min ,
  \qquad
  S = [0,1] \times [0,1] .
\]
Множество $S$ — выпуклый компакт, $\varphi$ выпукла и дифференцируема,
\[
  \nabla\varphi(x) = \begin{pmatrix} 2(x_1 - 0{,}3) \\ 2(x_2 - 0{,}7) \end{pmatrix} .
\]
Старт: $x_0 = (1,1)\T \in S$, $\varphi(x_0) = 0{,}49 + 0{,}09 = 0{,}58$.

\medskip
\noindent\textbf{Итерация 0.}
\[
  \nabla\varphi(x_0) = \begin{pmatrix} 1{,}4 \\ 0{,}6 \end{pmatrix} .
\]
Вспомогательная задача:
\[
  1{,}4\,x_1 + 0{,}6\,x_2 \ \longrightarrow\ \min_{x \in S} .
\]
Оба коэффициента положительны, поэтому минимум достигается в вершине с
наименьшими координатами:
\[
  y_0 = \begin{pmatrix} 0 \\ 0 \end{pmatrix} ,
  \qquad
  s_0 = y_0 - x_0 = \begin{pmatrix} -1 \\ -1 \end{pmatrix} ,
  \qquad
  \eta_0 = \nabla\T\varphi(x_0)\,s_0 = -1{,}4 - 0{,}6 = -2 < 0 .
\]
Шаг:
\[
  \varphi(x_0 + \alpha s_0)
   = (0{,}7 - \alpha)^2 + (0{,}3 - \alpha)^2
   = 0{,}58 - 2\alpha + 2\alpha^2 ,
\]
\[
  -2 + 4\alpha = 0 \quad\Longrightarrow\quad \alpha_0 = 0{,}5 \in [0,1] ,
\]
\[
  x_1 = (1-\alpha_0)x_0 + \alpha_0 y_0
      = \begin{pmatrix} 0{,}5 \\ 0{,}5 \end{pmatrix} ,
  \qquad
  \varphi(x_1) = 0{,}04 + 0{,}04 = 0{,}08 .
\]

\medskip
\noindent\textbf{Итерация 1.}
\[
  \nabla\varphi(x_1) = \begin{pmatrix} 0{,}4 \\ -0{,}4 \end{pmatrix} ,
  \qquad
  0{,}4\,x_1 - 0{,}4\,x_2 \ \longrightarrow\ \min_{x \in S} .
\]
Коэффициент при $x_1$ положителен — берём $x_1 = 0$; при $x_2$ отрицателен —
берём $x_2 = 1$:
\[
  y_1 = \begin{pmatrix} 0 \\ 1 \end{pmatrix} ,
  \qquad
  s_1 = \begin{pmatrix} -0{,}5 \\ 0{,}5 \end{pmatrix} ,
  \qquad
  \eta_1 = -0{,}2 - 0{,}2 = -0{,}4 < 0 .
\]
Шаг:
\[
  \varphi(x_1 + \alpha s_1)
   = (0{,}2 - 0{,}5\alpha)^2 + (-0{,}2 + 0{,}5\alpha)^2
   = 2\,(0{,}2 - 0{,}5\alpha)^2 ,
\]
\[
  \alpha_1 = 0{,}4 \in [0,1] ,
  \qquad
  x_2 = \begin{pmatrix} 0{,}5 - 0{,}2 \\ 0{,}5 + 0{,}2 \end{pmatrix}
      = \begin{pmatrix} 0{,}3 \\ 0{,}7 \end{pmatrix} = x_* ,
  \qquad \varphi(x_2) = 0 .
\]
Значения: $0{,}58 \to 0{,}08 \to 0$.
\end{example}

\begin{idea}
Три вещи, которые этот счёт иллюстрирует.

\begin{itemize}[nosep,left=0pt]
  \item Точки $y_k$ — всегда вершины квадрата: минимум линейной функции на
        многограннике достигается в вершине. Идти в $y_k$ целиком нельзя
        (при $\alpha = 1$ из $x_0$ мы попали бы в $(0,0)$ со значением
        $0{,}58$ — ровно столько же, сколько было), поэтому шаг обязателен.
  \item Допустимость получается даром: $x_{k+1} = (1-\alpha_k)x_k + \alpha_k y_k$
        — выпуклая комбинация двух точек $S$ при $\alpha_k \in [0,1]$. Никакого
        проектирования, в отличие от метода Розена.
  \item Попадание точно в $x_*$ за два шага — следствие того, что $x_*$ лежит
        \emph{внутри} $S$ и функция круговая. В типичной ситуации, когда
        оптимум на границе, метод подходит к нему сублинейно, с оценкой
        $O(1/k)$ из замечания выше: направления $s_k$ всё время «дёргаются»
        между вершинами.
\end{itemize}
\end{idea}
